(1) 两主轴传播相位分别为 \(\phi_x=2\pi n_xd/\lambda\)、\(\phi_y=2\pi n_yd/\lambda\)。剔除公共相位 \(e^{i\phi_x}\) 后，相对延迟为 \(\delta=\phi_y-\phi_x=2\pi(n_y-n_x)d/\lambda\)，Jones 矩阵可写 \(J=\operatorname{diag}(1,e^{i\delta})\)。无损波片的矩阵满足 \(J^\dagger J=I\)，故不改变总光强。
(2) 主轴各有相同实幅度的入射 Jones 向量 \(\boldsymbol E_{\rm in}=(1,1)^T/\sqrt2\)。透射为 \(\boldsymbol E_{\rm out}=J\boldsymbol E_{\rm in}=(1,e^{i\delta})^T/\sqrt2\)，其范数平方 \(\boldsymbol E_{\rm out}^\dagger\boldsymbol E_{\rm out}=(1+1)/2=1\)。两分量的相对相位决定偏振椭圆形状，整体相位无可测影响。
(3) 圆偏振要求两正交分量等幅且相位差 \(\delta=\pm\pi/2\)（模 \(2\pi\)），所以最薄正厚度为 \(d_{\rm QWP}=\lambda/[4|n_y-n_x|]\)，圆偏振左右旋向由相位符号决定。与初始 \(45^\circ\) 线偏振正交的输出是 \((1,-1)^T/\sqrt2\)，要求 \(\delta=\pi\)（模 \(2\pi\)），故最薄半波片厚度 \(d_{\rm HWP}=\lambda/[2|n_y-n_x|]\)。
(4) 分析器透射 Jones 单位矢量为 \(\boldsymbol e=(1,1)^T/\sqrt2\)，其透过幅度 \(\boldsymbol e^\dagger\boldsymbol E_{\rm out}=(1+e^{i\delta})/2\)。所以总强度比 \(I_{\rm out}/I_{\rm in}=|(1+e^{i\delta})/2|^2=\cos^2(\delta/2)\)。当 \(\delta=0\) 透过率为1；当 \(\delta=\pi/2\) 为1/2；当 \(\delta=\pi\) 为0，清楚体现由偏振干涉而非吸收产生的分析器消光。
